\chapter{Final Checklists \& Quick Reference}
\label{ch:checklists}

\section{The Master Implementation Checklist}

This is the consolidated, actionable checklist from the remote-jobs playbook, cross-referenced to the chapters in
this guide that cover each item in detail.

\subsection{Technical Positioning (Chapters \ref{ch:projects}--\ref{ch:projects-c})}
\begin{itemize}[nosep]
  \item[$\square$] Choose one technical specialization and define the business problems you can solve.
  \item[$\square$] Build or improve a genuinely useful project, using AI where appropriate.
  \item[$\square$] Deploy, test, monitor, and document the project's important failure cases.
  \item[$\square$] Prepare evidence of your architectural decisions, trade-offs, and measurable results.
\end{itemize}

\subsection{Professional Presence (Chapter \ref{ch:linkedin})}
\begin{itemize}[nosep]
  \item[$\square$] Rewrite your LinkedIn headline to communicate value and specialization.
  \item[$\square$] Publish technical explanations, lessons from failures, and credible project updates.
  \item[$\square$] Keep your resume and public project links consistent and easy to evaluate.
\end{itemize}

\subsection{Direct Outreach (Chapters \ref{ch:outreach-engineering}, \ref{ch:direct-outreach}, \ref{ch:beyond-postings})}
\begin{itemize}[nosep]
  \item[$\square$] Create a researched list of relevant funded startups.
  \item[$\square$] Identify founders or technical leaders whose companies match your capabilities.
  \item[$\square$] Write a personalized message tied to a real company problem.
  \item[$\square$] Offer a focused demonstration or prototype when it makes sense.
  \item[$\square$] Track outreach, replies, conversations, referrals, and outcomes.
\end{itemize}

\subsection{Consistency \& Review (Chapter \ref{ch:roadmap})}
\begin{itemize}[nosep]
  \item[$\square$] Reserve weekly time for projects, content, and targeted outreach.
  \item[$\square$] Review which messages and demonstrations generate meaningful conversations.
  \item[$\square$] Improve the process regularly rather than relying on message volume alone.
\end{itemize}

\section{Quick-Reference Dashboard: Every Number in One Place}

Every target number, cadence, and rate mentioned across the four source documents, gathered in one table for fast
lookup.

\begin{longtable}{p{0.44\textwidth} p{0.2\textwidth} p{0.26\textwidth}}
\toprule
\textbf{Metric} & \textbf{Target / figure} & \textbf{Source chapter} \\
\midrule
\endhead
LinkedIn connection-acceptance rate & 40\%+ & \ref{ch:linkedin} \\
LinkedIn message response rate & 20\%+ & \ref{ch:linkedin} \\
Daily new connections & 25/day & \ref{ch:linkedin} \\
Daily substantive comments & 10+/day & \ref{ch:linkedin} \\
LinkedIn post frequency & 2$\times$/week minimum & \ref{ch:linkedin} \\
About section length & 150--250 words & \ref{ch:linkedin} \\
Daily outreach messages (Outreach Engineering) & 25--30/day & \ref{ch:outreach-engineering} \\
Follow-up cadence & every 2--3 days until yes/no & \ref{ch:outreach-engineering}, \ref{ch:direct-outreach} \\
Illustrative acceptance $\rightarrow$ reply $\rightarrow$ positive $\rightarrow$ call $\rightarrow$ convert
rates & 20\% / 20\% / 25\% / 50\% / 20\% & \ref{ch:outreach-engineering} \\
Illustrative outreach-to-hire funnel & 800 $\rightarrow$ 40 $\rightarrow$ 4 $\rightarrow$ 1 (0.125\% overall) &
\ref{ch:direct-outreach} \\
Prototype turnaround offered to founders & $\sim$48 hours & \ref{ch:direct-outreach} \\
Project scope (time to complete) & 2--4 weeks per project & \ref{ch:projects} \\
Project explainability & understandable by a recruiter in $\le$2 minutes & \ref{ch:projects} \\
README read time for the hook section & $\sim$30 seconds & \ref{ch:projects-c} \\
Demo load time & $<$5 seconds & \ref{ch:projects-c} \\
Cited remote salary range (US/EU funded startups) & \$60,000--\$100,000/year & \ref{ch:foundations} \\
Cited target-company funding range & \$5M--\$20M+ & \ref{ch:foundations} \\
Cited target-company size & 20--40 employees & \ref{ch:foundations} \\
Implementation roadmap length & 16 weeks, 7 phases & \ref{ch:roadmap} \\
Open-source PR target & 3 merged PRs in one recognizable repo & \ref{ch:open-source} \\
YC funnel (emails $\to$ opens $\to$ replies $\to$ demos $\to$ customer) & 800 $\to$ 400 $\to$ 40 $\to$ 10 $\to$ 1
& \ref{ch:cold-email-principles} \\
Cited average cold-email reply rate (2026) & $\approx$3.4\% & \ref{ch:open-source} \\
Warm-introduction conversion multiplier (cited) & 2--3$\times$ ordinary cold email & \ref{ch:cold-email-principles}
\\
Professor cold-email length & 150--180 words & \ref{ch:cold-email-templates} \\
Professor outreach batch size & 7--10 professors, one per target group & \ref{ch:cold-email-templates} \\
Recruiter cold-email send time (cited) & Tue, 8--10\,AM local & \ref{ch:cold-email-templates} \\
1,200-email case study results & 60.2\% open / 10.2\% reply / 9$\to$12 opportunities & \ref{ch:case-study-emails}
\\
Cold-email sending ramp-up (case study) & $\sim$80/day (4 mailboxes) scaling to 40/mailbox & \ref{ch:case-study-emails}
\\
\bottomrule
\end{longtable}

\begin{caution}
\textbf{This guide now contains four separate, non-reconcilable funnel/rate models} — Puneet's funnel
($1{,}250\to215\to65\to15\to2$, Chapter~\ref{ch:outreach-engineering}), the playbook's funnel
($800\to40\to4\to1$, Chapter~\ref{ch:direct-outreach}), the YC funnel ($800\to400\to40\to10\to1$,
Chapter~\ref{ch:cold-email-principles}), and the 1,200-email case study's reported rates
(Chapter~\ref{ch:case-study-emails}). Each comes from a different source, context, and audience. Do not average
them or treat any one as predictive of your own results — track your own numbers and use whichever model's
\emph{structure} (not specific rates) helps you diagnose where your own funnel is weakest.
\end{caution}

\begin{caution}
As emphasized throughout this guide, the funnel and salary figures above are illustrative claims from the source
material, not guarantees. The LinkedIn engagement targets (40\% acceptance, 20\% response) are reasonable
benchmarks to instrument your own profile against, per the toolkit's own ``treat your profile like a production
system'' framing in Chapter~\ref{ch:linkedin}.
\end{caution}

